% Diapositivas de la exposición del miniproyecto de Visión Artificial (30 nov - 1 dic de 2026).
% Mismo estilo que la presentación de Machine Learning (plantilla de Pablo).
% Compilar desde esta carpeta con: pdflatex presentacion.tex  (dos veces).
% En Overleaf: subir este .tex y la carpeta figuras/ tal cual.
% Las imágenes están en figuras/ (si falta alguna, se busca en ../results/figuras/):
%   statefarm_clases.png y mrl_ojos.png -> uv run python scripts/figuras_dataset.py statefarm|mrl
%   dmd_camaras.png                     -> figura oficial del DMD (Vicomtech), sin modificar
%   demo.png                            -> captura de la demo (uv run python scripts/demo_webcam.py)
% Si no están, sale un recuadro en su lugar.
\documentclass[aspectratio=169,9pt,t,xcolor={table}]{beamer}

\usepackage[T1]{fontenc}
\usepackage[utf8]{inputenc}
\usepackage[spanish,es-noshorthands,es-tabla]{babel}
% Títulos en Oswald, una fuente estrecha. Si no está instalada (en Overleaf no lo está), se usa
% Helvetica Narrow, también estrecha, para que los títulos ocupen lo mismo.
\IfFileExists{oswald.sty}
  {\usepackage{oswald}\newcommand{\titular}{\sffamily}}
  {\renewcommand{\sfdefault}{phv}\newcommand{\titular}{\sffamily\fontseries{mc}\selectfont}}
\usepackage[default]{sourcecodepro} % texto general monoespaciado
\usepackage{tikz}
\usetikzlibrary{positioning,fit,calc,arrows.meta}
\usepackage[most]{tcolorbox}
\usepackage{array}

% ---------- Colores ----------
\definecolor{azul}{RGB}{61,120,186}
\definecolor{celeste}{HTML}{DAE6F3}
\definecolor{titulo}{HTML}{434343}
\definecolor{texto}{HTML}{333333}
\definecolor{marino}{HTML}{0B3954}
\definecolor{gris}{HTML}{7A7A7A}
\definecolor{lineas}{HTML}{A6A6A6}
\definecolor{cajafondo}{HTML}{FFF2CC}
\definecolor{cajaborde}{HTML}{BF9000}
\definecolor{cajatitulo}{HTML}{7F6000}
\definecolor{rojo}{HTML}{B23A3A}
\definecolor{verde}{HTML}{2E7D4F}

% ---------- Tema ----------
\renewcommand{\familydefault}{\ttdefault}
\AtBeginDocument{\renewcommand{\familydefault}{\ttdefault}\normalfont}
\setbeamertemplate{navigation symbols}{}
\setbeamersize{text margin left=11mm, text margin right=10mm}
\setbeamercolor{normal text}{fg=texto,bg=white}
\setbeamercolor{structure}{fg=azul}
\setbeamertemplate{frametitle}{%
  \nointerlineskip\vspace*{6mm}%
  {\titular\fontsize{20}{24}\selectfont\color{titulo}\insertframetitle}\par\vspace*{3mm}}
\setbeamertemplate{footline}{%
  \hfill{\ttfamily\fontsize{6}{6}\selectfont\color{gris}\insertframenumber}\hspace*{4mm}\vspace*{3mm}}
\setbeamertemplate{itemize item}{\color{azul}\textendash}
\setbeamertemplate{itemize subitem}{\color{azul}\textendash}
\setbeamertemplate{enumerate item}{\color{azul}\insertenumlabel.}
\setbeamerfont{itemize/enumerate body}{size=\normalsize}
\setbeamerfont{itemize/enumerate subbody}{size=\normalsize}
\setlength{\leftmargini}{5mm}

% Separador de sección: banda azul horizontal con el título en blanco
\newcommand{\seccion}[1]{%
  \begin{frame}
  \begin{tikzpicture}[remember picture,overlay]
    \fill[azul] ([yshift=27.4mm]current page.south west) rectangle ([yshift=62.6mm]current page.south east);
    \node[text=white,font=\titular\fontsize{38}{42}\selectfont] at (current page.center) {#1};
  \end{tikzpicture}
  \end{frame}}

% Caja amarilla con título
\newtcolorbox{caja}[2][]{enhanced, colback=cajafondo, colframe=cajaborde, boxrule=0.6pt,
  arc=5mm, left=4mm, right=3mm, top=2.5mm, bottom=2.5mm, #1,
  before upper={{\bfseries\color{cajatitulo}#2}\par\vspace{0.6mm}}}

% Caja celeste para los mensajes clave
\newtcolorbox{clave}[1][]{enhanced, colback=celeste, colframe=azul, boxrule=0.6pt, arc=3mm,
  left=3mm, right=3mm, top=1.5mm, bottom=1.5mm, #1}

% Imagen si existe; si no, un recuadro con el texto dado.
% Se busca primero en figuras/ (junto a este .tex, también en Overleaf) y después en
% ../results/figuras/ (donde las dejan los scripts del repositorio).
\newcommand{\recuadro}[3]{%
    \begin{tcolorbox}[enhanced, width=#1, height=#2, colback=white, colframe=lineas, boxrule=0.6pt,
      arc=2mm, borderline={0.6pt}{0pt}{lineas,dashed}, valign=center, halign=center]
      {\scriptsize\color{gris}#3}
    \end{tcolorbox}}
\newcommand{\figura}[4]{% nombre del fichero, ancho, alto, texto del recuadro
  \IfFileExists{figuras/#1}{\includegraphics[width=#2,height=#3,keepaspectratio]{figuras/#1}}{%
  \IfFileExists{../results/figuras/#1}{\includegraphics[width=#2,height=#3,keepaspectratio]{../results/figuras/#1}}{%
  \recuadro{#2}{#3}{#4}}}}

% Tablas con rejilla gris
\renewcommand{\arraystretch}{1.45}
\arrayrulecolor{lineas}

\hypersetup{pdftitle={Monitor del conductor},
  pdfauthor={Enzo Castro, Rodrigo Martínez y Pablo Sánchez}}

\begin{document}
\color{texto}

% =====================================================================
% Portada
% =====================================================================
\begin{frame}[plain]
\begin{tikzpicture}[remember picture,overlay]
  \fill[azul] (current page.north west) rectangle ([yshift=35.3mm]current page.south east);
  \fill[azul] ([xshift=-3.4mm,yshift=35.3mm]current page.south) -- ([xshift=3.4mm,yshift=35.3mm]current page.south)
              -- ([yshift=31.8mm]current page.south) -- cycle;
  \node[text=white,font=\titular\fontsize{44}{48}\selectfont] at ([yshift=-30mm]current page.north)
    {Monitor del conductor};
  \node[text=white,font=\ttfamily\normalsize] at ([yshift=-44mm]current page.north)
    {Distracción y somnolencia con cámaras interiores};
  \node[align=center,font=\ttfamily\normalsize,text=marino] at ([yshift=24mm]current page.south)
    {Enzo Castro · Rodrigo Martínez · Pablo Sánchez};
  \node[align=center,font=\ttfamily\footnotesize,text=texto] at ([yshift=16mm]current page.south)
    {Miniproyecto · Visión Artificial\\[0.6mm]Máster en Inteligencia Artificial · Universidad de Murcia};
  \node[font=\ttfamily\footnotesize,text=gris] at ([yshift=7mm]current page.south) {Noviembre de 2026};
\end{tikzpicture}
\end{frame}

\begin{frame}{Índice}
\vspace{6mm}
\begin{itemize}\setlength{\itemsep}{5mm}\large
  \item El problema y los datos
  \item El enfoque: tres señales y una decisión
  \item Cómo lo evaluaremos y qué tenemos ya
\end{itemize}
\end{frame}

% =====================================================================
% 1. El problema y los datos
% =====================================================================
\seccion{El problema y los datos}

\begin{frame}{El problema}
\vspace{1mm}
\begin{columns}[T,onlytextwidth]
\column{0.48\textwidth}
\textbf{Distracción: el factor n.º 1}
\vspace{1.5mm}
\begin{itemize}\setlength{\itemsep}{1.4mm}
  \item Presente en el \textbf{32 \%} de los siniestros mortales de 2025 en España (444).
  \item El \textbf{47,4 \%} de las denuncias de la última campaña de la DGT, por el móvil.
\end{itemize}
\column{0.48\textwidth}
\textbf{Ya es obligatorio en la UE}
\vspace{1.5mm}
\begin{itemize}\setlength{\itemsep}{1.4mm}
  \item Aviso de \textbf{somnolencia} (DDAW): todos los coches nuevos desde julio de 2024.
  \item Aviso de \textbf{distracción} (ADDW): desde el 7 de julio de 2026.
\end{itemize}
\end{columns}

\vspace{6mm}
\begin{clave}
\color{marino}\textbf{Objetivo:} un prototipo que, con cámaras dentro del coche, detecte si el
conductor está distraído o tiene sueño y le avise.
\end{clave}

\vspace{3mm}
{\scriptsize\color{gris}Fuentes: nota de prensa de la DGT (5 de octubre de 2026); Reglamento (UE) 2019/2144.}
\end{frame}

\begin{frame}{Los datos}
\centering\footnotesize
\setlength{\tabcolsep}{4pt}
\vspace{1mm}
\begin{tabular}{|p{30mm}|p{52mm}|p{48mm}|}
\hline
\textbf{Dataset} & \textbf{Qué contiene} & \textbf{Para qué} \\ \hline
\textbf{State Farm} (Kaggle) & 22.424 imágenes de 640$\times$480, 10 clases, 26 conductores, vista lateral &
  Entrenar y evaluar qué hace con las manos \\ \hline
\textbf{MRL Eye} & 84.898 recortes de ojos en infrarrojo, abiertos o cerrados, 37 sujetos &
  Ojo abierto o cerrado \\ \hline
\textbf{DMD} (Vicomtech) & Vídeo RGB de tres cámaras sincronizadas (cara, cuerpo, manos), etiquetas por fotograma &
  Mirada, somnolencia y el sistema completo \\ \hline
\textbf{Grabaciones propias} & Webcam o coche aparcado, solo nosotros &
  Probar fuera de los datasets \\ \hline
\end{tabular}

\vspace{4mm}
\raggedright\normalsize
\begin{itemize}\setlength{\itemsep}{1.2mm}
  \item El test oficial de State Farm \textbf{no tiene etiquetas}: train, validación y test salen
        de las imágenes etiquetadas.
  \item Las clases distinguen la mano: «escribe con la derecha» y «con la izquierda» son clases distintas.
\end{itemize}
\end{frame}

\begin{frame}{Los datos en imágenes}
\vspace{-1mm}
\centering
\figura{statefarm_clases.png}{\textwidth}{47mm}{State Farm: una imagen de cada clase (falta la figura)}

\vspace{1mm}
\figura{mrl_ojos.png}{\textwidth}{24mm}{MRL Eye: ojos abiertos y cerrados (falta la figura)}
\end{frame}

\begin{frame}{Una cámara no basta}
\vspace{1mm}
\centering
\figura{dmd_camaras.png}{\textwidth}{27mm}{DMD: cámaras de cara, cuerpo y manos (falta la figura)}

{\tiny\color{gris}Cámaras de cara, cuerpo y manos del DMD. Imagen: Vicomtech (CC BY-NC-ND 4.0).}

\vspace{2mm}
\raggedright
\begin{columns}[T,onlytextwidth]
\column{0.56\textwidth}
\footnotesize
\setlength{\tabcolsep}{4pt}
\renewcommand{\arraystretch}{1.3}
\begin{tabular}{|p{42mm}|p{24mm}|}
\hline
\textbf{¿MediaPipe ve la cara?} & \textbf{Detectada} \\ \hline
State Farm, de lado (una foto por clase) & 4 de 10 \\ \hline
DMD, cámara del cuerpo (de lado) & 41 de 120 \\ \hline
DMD, cámara de la cara (de frente) & 116 de 120 \\ \hline
\end{tabular}
\column{0.40\textwidth}
\small
De lado solo ve la cara cuando el conductor \textbf{deja de mirar la carretera}.

\vspace{2mm}
\begin{clave}
\color{marino}Como en los coches reales: una cámara de frente para ojos y mirada, y otra de lado
para las manos.
\end{clave}
\end{columns}
\end{frame}

\begin{frame}{Cómo separamos los datos}
\vspace{1mm}
Cada conductor aparece en cientos de fotogramas casi iguales. Si se reparten al azar, el mismo
conductor cae en train y en test.

\vspace{1mm}
\centering
\begin{tikzpicture}[x=1mm,y=0.8mm,
    caj/.style={draw=azul, rounded corners=2mm, minimum width=44mm, minimum height=9mm},
    c/.style={circle, minimum size=5.5mm, inner sep=0pt, font=\ttfamily\scriptsize\bfseries, text=white}]
  % Partición aleatoria
  \node[font=\ttfamily\small\bfseries, text=rojo] at (24,22) {Aleatoria};
  \node[caj, fill=celeste] (tr1) at (24,12) {};
  \node[caj, fill=cajafondo, draw=cajaborde] (te1) at (24,-2) {};
  \node[font=\ttfamily\scriptsize, text=gris, anchor=east] at (1,12) {train};
  \node[font=\ttfamily\scriptsize, text=gris, anchor=east] at (1,-2) {test};
  \foreach \l/\x/\col in {A/9/azul,B/19/verde,C/29/rojo,D/39/cajatitulo} {
    \node[c, fill=\col] at (\x,12) {\l};
    \node[c, fill=\col] at (\x,-2) {\l};
  }
  % Partición por conductor
  \node[font=\ttfamily\small\bfseries, text=verde] at (100,22) {Por conductor};
  \node[caj, fill=celeste] at (100,12) {};
  \node[caj, fill=cajafondo, draw=cajaborde] at (100,-2) {};
  \node[font=\ttfamily\scriptsize, text=gris, anchor=east] at (77,12) {train};
  \node[font=\ttfamily\scriptsize, text=gris, anchor=east] at (77,-2) {test};
  \foreach \l/\x/\col in {A/90/azul,B/100/verde,C/110/rojo} { \node[c, fill=\col] at (\x,12) {\l}; }
  \node[c, fill=cajatitulo] at (100,-2) {D};
  \node[font=\ttfamily\scriptsize, text=rojo] at (24,-12) {el modelo reconoce a la persona};
  \node[font=\ttfamily\scriptsize, text=verde] at (100,-12) {conductor nuevo, como en un coche real};
\end{tikzpicture}

\vspace{2mm}
\begin{clave}
\color{marino}Separamos \textbf{por conductor} (unos 18 / 4 / 4 en State Farm, también en el DMD) y por
sujeto en MRL Eye. Con la partición aleatoria, el acierto sale inflado: lo mediremos.
\end{clave}
\end{frame}

% =====================================================================
% 2. El enfoque
% =====================================================================
\seccion{El enfoque}

\begin{frame}{Tres señales y una decisión}
\vspace{1mm}
\centering
\begin{tikzpicture}[x=1mm,y=0.9mm,
    nodo/.style={draw=azul, fill=celeste, rounded corners=3mm, align=center, text=marino,
                 font=\ttfamily\scriptsize, minimum height=9mm, text width=24mm},
    flecha/.style={-{Stealth[length=2mm]}, draw=gris, line width=0.7pt},
    bt/.style={font=\ttfamily\tiny, text=cajatitulo}]
  \node[nodo, text width=15mm] (camb) at (0,11) {\textbf{Cámara}\\de lado};
  \node[nodo, text width=15mm] (camf) at (0,-11) {\textbf{Cámara}\\de frente};
  \node[nodo] (img)  at (29,11) {\textbf{Imagen completa}\\preprocesada};
  \node[nodo] (mp)   at (29,-11) {\textbf{Puntos de la cara}\\MediaPipe (478)};
  \node[nodo] (act)  at (61,18) {\textbf{1. Manos}\\qué está haciendo};
  \node[nodo] (ojos) at (61,0)  {\textbf{2. Ojos}\\abiertos o cerrados};
  \node[nodo] (mir)  at (61,-18) {\textbf{3. Mirada}\\giro de la cabeza};
  \node[nodo, fill=cajafondo, draw=cajaborde, text=cajatitulo, text width=20mm] (dec) at (93,0) {\textbf{Decisión}\\\textbf{temporal}};
  \node[nodo, fill=azul, text=white, text width=13mm] (alerta) at (118,0) {\textbf{Alerta}};
  \draw[flecha] (camb) -- (img);
  \draw[flecha] (camf) -- (mp);
  \draw[flecha] (img) -- (act);
  \draw[flecha] (mp) -- (ojos);
  \draw[flecha] (mp) -- (mir);
  \draw[flecha] (act) -- (dec);
  \draw[flecha] (ojos) -- (dec);
  \draw[flecha] (mir) -- (dec);
  \draw[flecha] (dec) -- (alerta);
  \node[bt, anchor=south] at (act.north) {CNN · ResNet · ViT · DINOv2/v3};
  \node[bt, anchor=north] at (ojos.south) {EAR · HOG+SVM · CNN};
  \node[bt, anchor=north] at (mir.south) {PnP + RANSAC};
  \node[bt, anchor=north] at (dec.south) {PERCLOS · tiempos};
\end{tikzpicture}

\vspace{2mm}
\raggedright\small
Un solo sistema para los dos avisos de la UE: \textbf{somnolencia} (ojos) y
\textbf{distracción} (manos y mirada).
\end{frame}

\begin{frame}{1. Manos: qué está haciendo}
\vspace{3mm}
\begin{columns}[T,onlytextwidth]
\column{0.47\textwidth}
\begin{caja}{Preprocesado}
\small
\begin{itemize}\setlength{\itemsep}{0.6mm}
  \item Redimensionar manteniendo 4:3.
  \item Normalizar como ImageNet.
  \item \textit{Augmentation}: recorte y giro leves, brillo y contraste.
  \item \textbf{Sin volteo horizontal}: «móvil con la derecha» pasaría a «con la izquierda».
\end{itemize}
\end{caja}
\column{0.49\textwidth}
\centering\footnotesize
\setlength{\tabcolsep}{4pt}
\begin{tabular}{|p{33mm}|p{30mm}|}
\hline
\textbf{Modelo} & \textbf{Qué aporta} \\ \hline
\textbf{CNN propia} & Línea base desde cero \\ \hline
\textbf{ResNet} & Congelada frente a \textit{fine-tuning} \\ \hline
\textbf{ViT} preentrenado & Atención frente a convolución \\ \hline
\textbf{DINOv2/v3} + SVM & Embeddings sin reentrenar \\ \hline
\textbf{Grad-CAM} & ¿Mira las manos o el fondo? \\ \hline
\end{tabular}
\end{columns}
\end{frame}

\begin{frame}{2. Ojos y 3. Mirada}
\vspace{1mm}
\begin{columns}[T,onlytextwidth]
\column{0.48\textwidth}
\textbf{Ojos: ¿abiertos o cerrados?}
\vspace{1.5mm}
\begin{caja}{EAR: alto del ojo entre ancho}
\centering
$\mathrm{EAR} = \dfrac{\lVert p_2-p_6\rVert + \lVert p_3-p_5\rVert}{2\,\lVert p_1-p_4\rVert}$

\vspace{1mm}
\raggedright\scriptsize Baja casi a 0 al cerrar el ojo. Seis puntos de MediaPipe por ojo.
\end{caja}
\vspace{1.5mm}
\begin{itemize}\setlength{\itemsep}{1mm}\small
  \item Geométrico (EAR) frente a \textbf{HOG + SVM} y una \textbf{CNN} sobre MRL Eye.
  \item Recortes: escala de grises y \textbf{CLAHE} para poca luz e infrarrojo.
\end{itemize}
\column{0.48\textwidth}
\textbf{Mirada: ¿hacia dónde gira la cabeza?}
\vspace{1.5mm}
\begin{itemize}\setlength{\itemsep}{1.4mm}\small
  \item \textbf{PnP}: de los puntos 2D de la cara y un modelo 3D de cara a la rotación de la cabeza.
  \item \textbf{RANSAC} frente a mínimos cuadrados: cuando la mano o el móvil tapan parte de la
        cara, hay puntos erróneos.
  \item Pose «mirando a la carretera» como referencia: rotación relativa y ángulos de Euler, como en
        la entrega del Tema 2. Da igual dónde esté la cámara.
\end{itemize}
\end{columns}
\end{frame}

\begin{frame}{Decisión temporal}
\vspace{1mm}
Un fotograma no basta: un parpadeo cierra los ojos y mirar el retrovisor gira la cabeza.

\vspace{3mm}
\centering\footnotesize
\begin{tabular}{|l|p{62mm}|l|}
\hline
\textbf{Alerta} & \textbf{Regla} & \textbf{Aviso UE} \\ \hline
\textbf{Microsueño} & Ojos cerrados seguidos más de 1 s & DDAW \\ \hline
\textbf{Somnolencia} & PERCLOS (\% del tiempo con los ojos cerrados en 30 s) por encima del umbral & DDAW \\ \hline
\textbf{Mirada fuera} & Cabeza girada más de 30° durante más de 2 s & ADDW \\ \hline
\textbf{Actividad} & Móvil, bebida, radio... durante varios fotogramas seguidos & ADDW \\ \hline
\end{tabular}

\vspace{4mm}
\raggedright
{\small\color{gris}Los umbrales son de partida: los ajustaremos con nuestras grabaciones y los justificaremos en la memoria.}
\end{frame}

\begin{frame}{Tópicos de la asignatura}
\centering\footnotesize
\setlength{\tabcolsep}{4pt}
\renewcommand{\arraystretch}{1.35}
\vspace{0mm}
\begin{tabular}{|p{12mm}|p{104mm}|r|}
\hline
\textbf{Bloque} & \textbf{Dónde lo usamos} & \textbf{\%} \\ \hline
\textbf{BT1} & Particiones por conductor y semillas · gráficas · demo en directo & 15 \\ \hline
\textbf{BT2} & Giro de la cabeza con PnP & 5 \\ \hline
\textbf{BT3} & CLAHE para poca luz e infrarrojo & 5 \\ \hline
\textbf{BT4} & HOG + SVM para el ojo & 5 \\ \hline
\textbf{BT5} & CNN propia · transfer learning · augmentation · métricas · clásico frente a aprendido · Grad-CAM & 30 \\ \hline
\textbf{BT6} & ViT · DINOv2/v3 + SVM · puntos de la cara · varios modelos juntos & 25 \\ \hline
\textbf{BT7} & RANSAC frente a mínimos cuadrados · sensibilidad al umbral & 15 \\ \hline
\rowcolor{celeste}\textbf{Total} & & \textbf{100} \\ \hline
\end{tabular}

\vspace{3mm}
\raggedright\normalsize
Todo en \textbf{un único proyecto}: cada técnica está donde la pide el problema.
\end{frame}

% =====================================================================
% 3. Evaluación y estado
% =====================================================================
\seccion{Evaluación y estado}

\begin{frame}{Cómo lo evaluaremos}
\vspace{1mm}
\begin{columns}[T,onlytextwidth]
\column{0.48\textwidth}
\textbf{Métricas}
\vspace{1.5mm}
\begin{itemize}\setlength{\itemsep}{1.4mm}
  \item Accuracy, \textbf{F1 macro} y matriz de confusión (10 clases).
  \item Acierto de ojo abierto o cerrado, por sujeto.
  \item Mirada frente a la etiqueta «mira a la carretera» del DMD.
  \item Fotogramas por segundo en un portátil.
\end{itemize}
\column{0.48\textwidth}
\textbf{Comparaciones}
\vspace{1.5mm}
\begin{itemize}\setlength{\itemsep}{1.4mm}
  \item Mismo modelo: partición \textbf{aleatoria} frente a \textbf{por conductor}.
  \item CNN frente a ViT y DINOv2/v3; clásico frente a aprendido.
  \item Cuánto cae el acierto al pasar de State Farm \textbf{al DMD} (otro coche, otros
        conductores) y a \textbf{nuestras grabaciones}.
\end{itemize}
\end{columns}

\vspace{6mm}
\begin{clave}
\color{marino}La pregunta central: ¿cuánto acierta de verdad con un conductor que nunca ha visto?
\end{clave}
\end{frame}

\begin{frame}{Lo que tenemos ya}
\vspace{1mm}
\begin{columns}[T,onlytextwidth]
\column{0.52\textwidth}
\centering
\figura{demo.png}{\linewidth}{38mm}{Captura de la demo\\[1mm](pendiente)}
\column{0.44\textwidth}
\begin{caja}{Demo inicial}
\begin{itemize}\setlength{\itemsep}{0.6mm}
  \item Con la webcam y sobre el DMD (dos cámaras).
  \item Ojos cerrados más de 1~s: aviso.
  \item Mirada fuera más de 2~s: aviso (giro de MediaPipe, por ahora).
\end{itemize}
\end{caja}
\end{columns}

\vspace{3mm}
\begin{clave}
\color{marino}\textbf{Primer número} (DMD, un conductor): nuestra «mirada fuera» coincide con su
etiqueta en el \textbf{92\,\%} de los fotogramas, pero solo detecta el \textbf{65\,\%} de las veces que
no mira la carretera.
\end{clave}
\end{frame}

\begin{frame}{Riesgos y plan B}
\vspace{2mm}
\begin{columns}[T,onlytextwidth]
\column{0.48\textwidth}
\textbf{\color{rojo}Riesgos}
\vspace{2mm}
\begin{itemize}\setlength{\itemsep}{1.8mm}
  \item El DMD es grande: unos 5 GB por conductor.
  \item Sus tres cámaras no empezaron a grabar a la vez.
  \item Las clases de State Farm y del DMD no coinciden del todo.
  \item Poca luz y gafas: los ojos se ven peor.
\end{itemize}
\column{0.48\textwidth}
\textbf{\color{azul}Plan B}
\vspace{2mm}
\begin{itemize}\setlength{\itemsep}{1.8mm}
  \item Pocos conductores y uno de cada 10-15 fotogramas.
  \item La herramienta DEx de Vicomtech las alinea.
  \item Comparar solo las clases comunes.
  \item CLAHE y comparar con el infrarrojo de MRL Eye.
\end{itemize}
\end{columns}
\end{frame}

\begin{frame}{Reparto y calendario}
\vspace{1mm}
\centering\footnotesize
\begin{tabular}{|l|p{90mm}|}
\hline
\textbf{Quién} & \textbf{Parte} \\ \hline
\textbf{Rodrigo} & Manos: CNN, ResNet y particiones por conductor \\ \hline
\textbf{Enzo} & Manos: ViT, DINOv2/v3 y Grad-CAM \\ \hline
\textbf{Pablo} & Ojos y mirada: EAR, HOG+SVM, PnP con RANSAC y la demo \\ \hline
\end{tabular}

\vspace{4mm}
\begin{tabular}{|l|p{90mm}|}
\hline
\textbf{Cuándo} & \textbf{Qué} \\ \hline
\textbf{2 – 13 dic} & Modelos principales y evaluación montada \\ \hline
\textbf{14 dic – 6 ene} & Comparativas, grabaciones propias, casos de fallo y memoria \\ \hline
\textbf{7 ene} & Entrega: memoria, código y vídeo \\ \hline
\end{tabular}
\end{frame}

% =====================================================================
% Cierre
% =====================================================================
\begin{frame}[plain]
\begin{tikzpicture}[remember picture,overlay]
  \fill[azul] (current page.north west) rectangle ([yshift=35.3mm]current page.south east);
  \fill[azul] ([xshift=-3.4mm,yshift=35.3mm]current page.south) -- ([xshift=3.4mm,yshift=35.3mm]current page.south)
              -- ([yshift=31.8mm]current page.south) -- cycle;
  \node[text=white,font=\titular\fontsize{48}{52}\selectfont] at ([yshift=-32mm]current page.north) {Gracias};
  \node[align=center,font=\ttfamily\normalsize,text=marino] at ([yshift=24mm]current page.south)
    {Enzo Castro · Rodrigo Martínez · Pablo Sánchez};
  \node[align=center,font=\ttfamily\footnotesize,text=texto] at ([yshift=16mm]current page.south)
    {Miniproyecto · Visión Artificial\\[0.6mm]Máster en Inteligencia Artificial · Universidad de Murcia};
  \node[font=\ttfamily\footnotesize,text=gris] at ([yshift=7mm]current page.south) {Noviembre de 2026};
\end{tikzpicture}
\end{frame}

\end{document}
